\chapter{Additional figures from the recordings}
\label{app:gallery}

\noindent\emph{Every figure here is drawn from data already described in
\cref{app:pilot,ch:vision,ch:results}; this chapter exists so that a reader
can see the whole record rather than the headline. Labels are in plain
words. Scripts: \repo{thesis\_v3/figures/gallery/*/make\_*\_gallery.py}.}

\clearpage
\section{The pilot, session by session}

\begin{figure}[H]
  \centering
  \includegraphics[width=\linewidth]{gallery/pilot/01_session_glance_P2.pdf}
  \caption[What one operator's sessions looked like, second by second]{One
  operator's five sessions (P2), second by second: shaded where the clutch
  was held and the operator was steering, the grey line the speed of the
  robot's hand, black where the emergency stop was held. The other four
  operators' strips are in the gallery directory.}
\end{figure}

\begin{figure}[H]
  \centering
  \begin{subfigure}[b]{0.49\linewidth}
    \centering\includegraphics[width=\linewidth]{gallery/pilot/02_paths_topdown.pdf}
    \caption{from above}
  \end{subfigure}\hfill
  \begin{subfigure}[b]{0.49\linewidth}
    \centering\includegraphics[width=\linewidth]{gallery/pilot/02_paths_side.pdf}
    \caption{from the side}
  \end{subfigure}
  \caption[Where the robot's hand went in every VR session]{Where the
  robot's hand went in every one of the twenty-three VR sessions, coloured
  from the start (dark) to the end (yellow) of each; one cell per operator
  and condition.}
\end{figure}

\begin{figure}[H]
  \centering
  \begin{subfigure}[b]{0.49\linewidth}
    \centering\includegraphics[width=\linewidth]{gallery/pilot/03_order_effects.pdf}
    \caption{practice}
  \end{subfigure}\hfill
  \begin{subfigure}[b]{0.49\linewidth}
    \centering\includegraphics[width=\linewidth]{gallery/pilot/07_tracking_loss.pdf}
    \caption{controller tracking loss}
  \end{subfigure}
  \caption[Practice effects and tracking loss]{Left: each operator's
  sessions in the order recorded; no consistent drift with practice, so
  order is not obviously driving the direct/assisted difference. Right: the
  headset lost sight of the controller in a median \SI{0}{\percent} of
  samples, worst \SI{8}{\percent}, clustered at session edges.}
\end{figure}

\begin{figure}[H]
  \centering
  \includegraphics[width=\linewidth]{gallery/pilot/04_lag.pdf}
  \caption[How far behind the operator the robot ran]{Robot lag behind the
  operator while steering: median \SI{1.6}{\milli\metre} direct,
  \SI{1.7}{\milli\metre} with assistance; past the \SI{30}{\milli\metre}
  capture gate in \SI{0.5}{\percent} and \SI{0.9}{\percent} of samples; lag
  grows roughly linearly with hand speed.}
\end{figure}

\begin{figure}[H]
  \centering
  \includegraphics[width=\linewidth]{gallery/pilot/05_clutch_behaviour.pdf}
  \caption[How operators used the clutch]{Clutch use per session:
  engagements per minute (median 6.6 under both conditions), how long each
  engagement was held (\SI{3.8}{\second} direct, \SI{4.3}{\second}
  assisted), the pause before re-engaging, and refused engagements (3.8 per
  session under both).}
\end{figure}

\begin{figure}[H]
  \centering
  \begin{subfigure}[b]{0.49\linewidth}
    \centering\includegraphics[width=\linewidth]{gallery/pilot/12_robot_vs_hand_travel.pdf}
    \caption{robot path per metre of hand motion}
  \end{subfigure}\hfill
  \begin{subfigure}[b]{0.49\linewidth}
    \centering\includegraphics[width=\linewidth]{gallery/pilot/12_robot_speed.pdf}
    \caption{robot hand speed while steering}
  \end{subfigure}
  \caption[Robot path against operator motion, and robot speed]{Left: the
  operator's controller covers \SIrange{10}{50}{\metre} per session (raw
  motion, jitter included) while the robot's hand covers
  \SIrange{0.2}{7}{\metre}; the ratio is the same under both conditions.
  Right: the robot's hand moves at similar speeds under both conditions
  (medians \SI{35}{} and \SI{45}{\milli\metre\per\second}) -- assistance
  lengthens the path, it does not speed the robot up.}
\end{figure}

\clearpage
\section{Sim-to-real tracking, every session and every joint}

\begin{figure}[H]
  \centering
  \includegraphics[width=0.9\linewidth]{gallery/pilot/08_sim_real_heatmap.pdf}
  \caption[Tracking error, sessions by joints]{How closely the real arm
  followed the simulated one, every session (rows) by joint (columns), both
  cohorts. Median per-joint error \SI{1.44}{\degree}; the three M1 rows and
  two P5 rows light up on several joints at once, the signature of a bad
  input channel rather than a bad joint; every P1--P4 session stays under
  \SI{4.6}{\degree} on every joint.}
  \label{fig:sim-real-heatmap}
\end{figure}

\begin{figure}[H]
  \centering
  \includegraphics[width=\linewidth]{gallery/pilot/08_sim_real_error.pdf}
  \caption[Tracking error per joint, and its growth with speed]{Pooled over
  the VR cohort: the error is centred on zero for every joint and grows with
  joint speed (median \SI{0.13}{\degree} at rest, \SI{2.59}{\degree} in the
  fastest bin) -- the real arm lags a moving command rather than sitting off
  a still one.}
\end{figure}

\begin{figure}[H]
  \centering
  \begin{subfigure}[b]{0.55\linewidth}
    \centering\includegraphics[width=\linewidth]{gallery/pilot/09_master_raw_channels.pdf}
    \caption{raw master channels, a good and a bad session}
  \end{subfigure}\hfill
  \begin{subfigure}[b]{0.43\linewidth}
    \centering\includegraphics[width=\linewidth]{gallery/pilot/09_master_hand_paths.pdf}
    \caption{the master hand's path, all nine sessions}
  \end{subfigure}
  \caption[The master arm's channels and hand paths]{Left: the seven raw
  sensor channels of one M2 session (tracking error \SI{2.2}{\degree})
  beside one M1 session (\SI{21.2}{\degree}); a channel that traces a
  trajectory is alive, one that sits flat or jumps between levels is the
  defect the health verdict names. Right: M1's hand paths are wide, jagged
  sweeps; M2's and M3's are compact.}
\end{figure}

\clearpage
\section{The scripted task set, in pictures}

\begin{figure}[H]
  \centering
  \includegraphics[width=\linewidth]{gallery/verification/05_circuit_box_by_mode.pdf}
  \caption[The circuit-box task in every mode]{The circuit-box-and-multimeter
  task under all five modes at four moments (simulation). Both arms lift the
  box and bring the meter to it in every mode; the runs differ in timing --
  grasps at \SIrange{4.4}{4.8}{\second} in four modes, \SI{37.0}{\second}
  through the VR path -- not in outcome.}
\end{figure}

\begin{figure}[H]
  \centering
  \includegraphics[width=\linewidth]{gallery/verification/06_reach_target_by_mode.pdf}
  \caption[The reach task in every mode]{Reaching three targets under all
  five modes, at the start, a third, two thirds and the end of the run
  (simulation). \SI{16.5}{\second} in four modes, \SI{49.1}{\second} through
  the VR path.}
\end{figure}

\begin{figure}[H]
  \centering
  \begin{subfigure}[b]{0.49\linewidth}
    \centering\includegraphics[width=\linewidth]{gallery/verification/03_gripper_traces_by_mode.pdf}
    \caption{finger angle through every grasping run}
  \end{subfigure}\hfill
  \begin{subfigure}[b]{0.49\linewidth}
    \centering\includegraphics[width=\linewidth]{gallery/verification/12_tray_tilt_by_mode.pdf}
    \caption{tray tilt while carrying}
  \end{subfigure}
  \caption[Gripper closure and tray tilt, by mode]{Left: in every mode the
  fingers close when the log says an object was grasped (\SI{24.2}{\degree}
  on a \SI{40}{\milli\metre} cube, \SI{29.7}{\degree} on the
  \SI{30}{\milli\metre} objects) and open when it says released. Right: the
  tray stays within a degree of level in four modes; the VR-direct run
  spiked to \SI{9.9}{\degree} for \SI{0.31}{\second}, past the
  \SI{6.8}{\degree} roll-off angle.}
\end{figure}

\begin{figure}[H]
  \centering
  \begin{subfigure}[b]{0.49\linewidth}
    \centering\includegraphics[width=\linewidth]{gallery/verification/08_placement_error_history.pdf}
    \caption{set-down error across four recorded versions}
  \end{subfigure}\hfill
  \begin{subfigure}[b]{0.49\linewidth}
    \centering\includegraphics[width=\linewidth]{gallery/verification/13_scripted_run_durations.pdf}
    \caption{run duration by mode}
  \end{subfigure}
  \caption[Placement error history and run durations]{Left: how far each
  cube was set down from its pad across the four recorded versions of the
  pick-and-place -- means \SIlist{130;174;312}{\milli\metre}, then
  \SI{56}{\milli\metre} on every cube. Right: the same scripted task takes
  about three times longer through the VR-controller path (\SI{49}{} against
  \SI{16.5}{\second} for the reach task), which is rate-limited by its own
  filter.}
\end{figure}

\begin{figure}[H]
  \centering
  \includegraphics[width=\linewidth]{gallery/verification/09_robot_travel_no_operator.pdf}
  \caption[Same waypoints, no operator: robot travel per mode]{With no
  person in the loop, the same commanded waypoints make the arm travel
  further under the assisted modes: reach task, left arm, \SI{1.41}{\metre}
  direct against \SI{2.03}{\metre} with mannequin-master assistance
  (\SI{+45}{\percent}); \SI{1.46}{} against \SI{1.52}{\metre} for the VR
  pair (\SI{+4}{\percent}). This is the no-operator baseline behind
  \cref{sec:pilot-limits}.}
  \label{fig:mode-travel}
\end{figure}

\begin{figure}[H]
  \centering
  \begin{subfigure}[b]{0.49\linewidth}
    \centering\includegraphics[width=\linewidth]{gallery/verification/10_vr_lag_vs_hand_speed.pdf}
    \caption{VR lag against hand speed}
  \end{subfigure}\hfill
  \begin{subfigure}[b]{0.49\linewidth}
    \centering\includegraphics[width=\linewidth]{gallery/verification/11_master_sweep_commanded_vs_moved.pdf}
    \caption{master sweep: asked against moved}
  \end{subfigure}
  \caption[VR lag against speed, and the master sweep]{Left: in a scripted
  VR protocol (simulation, 2026-08-26) the robot never fell more than
  \SI{22.7}{\milli\metre} behind a hand moving under
  \SI{0.27}{\metre\per\second}; a deliberately fast push at
  \SI{0.35}{\metre\per\second} left it \SIrange{83}{93}{\milli\metre} behind,
  three times the capture gate. Right: twenty-eight set movements of the
  mannequin master: the left arm moved roughly what was asked, the right
  undershot on eight of fourteen; the hand never came within half a metre of
  the wearer.}
\end{figure}

\begin{figure}[H]
  \centering
  \includegraphics[width=0.8\linewidth]{gallery/verification/07_seven_camera_angles.pdf}
  \caption[One grasp, seven cameras]{The first grasp of the autonomous
  pick-and-place from the seven fixed viewpoints plus the gripper's own
  camera; every run was filmed this way so a grasp is judged from the
  pictures rather than asserted.}
\end{figure}

\clearpage
\section{Every trajectory, every session}

\begin{figure}[H]
  \centering
  \includegraphics[width=0.47\linewidth]{gallery/trajectories/composite_best_per_operator.pdf}
  \caption[Each VR operator's cleanest direct and assisted session]{Each of
  the five VR operators' cleanest direct (left) and assisted (right) session:
  the operator's hand (faint), the simulated robot's hand, and the real
  robot's hand computed from its own joint encoders; the first clutch press
  and release are marked, and each panel label carries the typical gap
  between simulated and real hand.}
\end{figure}

\begin{figure}[H]
  \centering
  \includegraphics[width=\linewidth]{gallery/trajectories/3d_grid_P1.pdf}
  \caption[Every session of P1 in three dimensions]{Every recorded session
  of P1: the operator's hand (faint, smoothed), the simulated robot's hand
  and the real robot's hand from its own encoders.}
\end{figure}

\begin{figure}[H]
  \centering
  \includegraphics[width=\linewidth]{gallery/trajectories/3d_grid_P2.pdf}
  \caption[Every session of P2 in three dimensions]{Every recorded session
  of P2: the operator's hand (faint, smoothed), the simulated robot's hand
  and the real robot's hand from its own encoders.}
\end{figure}

\begin{figure}[H]
  \centering
  \includegraphics[width=\linewidth]{gallery/trajectories/3d_grid_P3.pdf}
  \caption[Every session of P3 in three dimensions]{Every recorded session
  of P3: the operator's hand (faint, smoothed), the simulated robot's hand
  and the real robot's hand from its own encoders.}
\end{figure}

\begin{figure}[H]
  \centering
  \includegraphics[width=\linewidth]{gallery/trajectories/3d_grid_P4.pdf}
  \caption[Every session of P4 in three dimensions]{Every recorded session
  of P4: the operator's hand (faint, smoothed), the simulated robot's hand
  and the real robot's hand from its own encoders.}
\end{figure}

\begin{figure}[H]
  \centering
  \includegraphics[width=\linewidth]{gallery/trajectories/3d_grid_P5.pdf}
  \caption[Every session of P5 in three dimensions]{Every recorded session
  of P5: the operator's hand (faint, smoothed), the simulated robot's hand
  and the real robot's hand from its own encoders.}
\end{figure}

\begin{figure}[H]
  \centering
  \includegraphics[width=\linewidth]{gallery/trajectories/3d_grid_M1.pdf}
  \caption[Every session of M1 in three dimensions]{Every recorded session
  of M1: the operator's hand (faint, smoothed), the simulated robot's hand
  and the real robot's hand from its own encoders.}
\end{figure}

\begin{figure}[H]
  \centering
  \includegraphics[width=\linewidth]{gallery/trajectories/3d_grid_M2.pdf}
  \caption[Every session of M2 in three dimensions]{Every recorded session
  of M2: the operator's hand (faint, smoothed), the simulated robot's hand
  and the real robot's hand from its own encoders.}
\end{figure}

\begin{figure}[H]
  \centering
  \includegraphics[width=\linewidth]{gallery/trajectories/3d_grid_M3.pdf}
  \caption[Every session of M3 in three dimensions]{Every recorded session
  of M3: the operator's hand (faint, smoothed), the simulated robot's hand
  and the real robot's hand from its own encoders.}
\end{figure}

\begin{figure}[H]
  \centering
  \includegraphics[width=\linewidth]{gallery/trajectories/gap_over_time_P1.pdf}
  \caption[Gap between simulated and real hand through each of P1's sessions]{The
  distance between the simulated and the real robot hand through each of
  P1's sessions, with the session mean in each label and the
  \SI{30}{\milli\metre} grasp gate drawn. A flat tail at the end of a
  session is the real-arm columns freezing while the end-of-session
  emergency stop is held, not a standing offset.}
\end{figure}

\begin{figure}[H]
  \centering
  \includegraphics[width=\linewidth]{gallery/trajectories/gap_over_time_P2.pdf}
  \caption[Gap between simulated and real hand through each of P2's sessions]{The
  distance between the simulated and the real robot hand through each of
  P2's sessions, with the session mean in each label and the
  \SI{30}{\milli\metre} grasp gate drawn. A flat tail at the end of a
  session is the real-arm columns freezing while the end-of-session
  emergency stop is held, not a standing offset.}
\end{figure}

\begin{figure}[H]
  \centering
  \includegraphics[width=\linewidth]{gallery/trajectories/gap_over_time_P3.pdf}
  \caption[Gap between simulated and real hand through each of P3's sessions]{The
  distance between the simulated and the real robot hand through each of
  P3's sessions, with the session mean in each label and the
  \SI{30}{\milli\metre} grasp gate drawn. A flat tail at the end of a
  session is the real-arm columns freezing while the end-of-session
  emergency stop is held, not a standing offset.}
\end{figure}

\begin{figure}[H]
  \centering
  \includegraphics[width=\linewidth]{gallery/trajectories/gap_over_time_P4.pdf}
  \caption[Gap between simulated and real hand through each of P4's sessions]{The
  distance between the simulated and the real robot hand through each of
  P4's sessions, with the session mean in each label and the
  \SI{30}{\milli\metre} grasp gate drawn. A flat tail at the end of a
  session is the real-arm columns freezing while the end-of-session
  emergency stop is held, not a standing offset.}
\end{figure}

\begin{figure}[H]
  \centering
  \includegraphics[width=\linewidth]{gallery/trajectories/gap_over_time_P5.pdf}
  \caption[Gap between simulated and real hand through each of P5's sessions]{The
  distance between the simulated and the real robot hand through each of
  P5's sessions, with the session mean in each label and the
  \SI{30}{\milli\metre} grasp gate drawn. A flat tail at the end of a
  session is the real-arm columns freezing while the end-of-session
  emergency stop is held, not a standing offset.}
\end{figure}

\begin{figure}[H]
  \centering
  \includegraphics[width=\linewidth]{gallery/trajectories/gap_over_time_M1.pdf}
  \caption[Gap between simulated and real hand through each of M1's sessions]{The
  distance between the simulated and the real robot hand through each of
  M1's sessions, with the session mean in each label and the
  \SI{30}{\milli\metre} grasp gate drawn. A flat tail at the end of a
  session is the real-arm columns freezing while the end-of-session
  emergency stop is held, not a standing offset.}
\end{figure}

\begin{figure}[H]
  \centering
  \includegraphics[width=\linewidth]{gallery/trajectories/gap_over_time_M2.pdf}
  \caption[Gap between simulated and real hand through each of M2's sessions]{The
  distance between the simulated and the real robot hand through each of
  M2's sessions, with the session mean in each label and the
  \SI{30}{\milli\metre} grasp gate drawn. A flat tail at the end of a
  session is the real-arm columns freezing while the end-of-session
  emergency stop is held, not a standing offset.}
\end{figure}

\begin{figure}[H]
  \centering
  \includegraphics[width=\linewidth]{gallery/trajectories/gap_over_time_M3.pdf}
  \caption[Gap between simulated and real hand through each of M3's sessions]{The
  distance between the simulated and the real robot hand through each of
  M3's sessions, with the session mean in each label and the
  \SI{30}{\milli\metre} grasp gate drawn. A flat tail at the end of a
  session is the real-arm columns freezing while the end-of-session
  emergency stop is held, not a standing offset.}
\end{figure}

Top and side views of every session, and each session's panel as its own
file, are in the trajectories gallery directory
(\repo{figures/gallery/} \repo{trajectories/}).

\clearpage
\section{The cameras}

\begin{figure}[H]
  \centering
  \includegraphics[width=\linewidth]{gallery/vision/pipeline_status.pdf}
  \caption[The 33 perception stages, per camera: result, refusal or not
  applicable]{Every stage of the picture-to-grasp pipeline on the four real
  cameras: 28, 29, 27 and 12 of 33 stages gave a result; every other outcome
  is a named refusal or a by-design absence, and none crashed.}
  \label{fig:pipeline-status}
\end{figure}

\begin{figure}[H]
  \centering
  \includegraphics[width=\linewidth]{gallery/vision/four_cameras.pdf}
  \caption[The four cameras' views]{The four cameras' views of the same
  moment: two ride on the hands and see the table; two stand across the room
  and see the whole rig.}
\end{figure}

\begin{figure}[H]
  \centering
  \includegraphics[width=\linewidth]{gallery/vision/plane_fit_per_camera.pdf}
  \caption[Where the flat table was found, per camera]{The support plane the
  pipeline fitted, recomputed over each camera's own depth array with the
  stored plane and tolerance: \SI{73.95}{\percent} of the left wrist
  camera's points lie on it, \SI{52.79}{\percent} of the right's, and under
  \SI{1}{\percent} of the room depth camera's (the pipeline's own log says
  \SI{3.75}{\percent} for that camera; both mean no usable surface).}
\end{figure}

\begin{figure}[H]
  \centering
  \includegraphics[width=0.8\linewidth]{gallery/vision/plan_view_left.pdf}
  \caption[What the left wrist camera believes is on the table]{The left
  wrist camera's point cloud from above, coloured by height, with the
  pipeline's oriented boxes: the box (red, too wide for the
  \SI{85}{\milli\metre} jaw) and the cube (green, graspable) plus three small
  fragments; nearest object \SI{0.539}{\metre}, furthest \SI{0.843}{\metre}.
  This is the map the planner plans against.}
\end{figure}

\begin{figure}[H]
  \centering
  \includegraphics[width=\linewidth]{gallery/vision/depth_coverage.pdf}
  \caption[How much of each picture has a depth reading]{Share of each
  picture with any depth reading (\SIlist{89;80;15}{\percent}) and within the
  working range (\SIlist{42;33;13}{\percent}); the room depth camera's
  returns lie mostly \SIrange{2}{6}{\metre} away -- a different instrument
  from the wrist cameras, not merely a distant one.}
\end{figure}

\begin{figure}[H]
  \centering
  \includegraphics[width=\linewidth]{gallery/vision/left_vs_right_pipeline.pdf}
  \caption[Both wrist cameras: picture, objects picked out, boxes]{The same
  scene from both hands: each wrist camera isolates the cube and the box.}
\end{figure}

\begin{figure}[H]
  \centering
  \includegraphics[width=\linewidth]{gallery/vision/arm_park_error.pdf}
  \caption[How far short each arm stopped]{The thirty-six calibration moves:
  every \SI{100}{\milli\metre} move stopped \SIrange{5.1}{8.3}{\milli\metre}
  short, the same for both arms and all three speeds, with the worst joint's
  rest error in a \SIrange{0.289}{0.296}{\degree} band (the pooled per-joint
  fit in \cref{sec:sim-to-real} gives \SI{0.305}{\degree}).}
\end{figure}

\clearpage
\section{Measured baselines, in pictures}

\begin{figure}[H]
  \centering
  \includegraphics[width=\linewidth]{gallery/baselines/error_budget.pdf}
  \caption[Where the positioning error comes from]{Each known source of
  finger-landing error as a bar, fixed offsets distinguished from random
  terms and each fix shown against the row it fixes; combined
  \SI{12.9}{\milli\metre} against the \SI{30}{\milli\metre} grasp line. The
  camera's true position on the wrist is unmeasured and not on this chart.}
\end{figure}

\begin{figure}[H]
  \centering
  \begin{subfigure}[b]{0.55\linewidth}
    \centering\includegraphics[width=\linewidth]{gallery/baselines/reaction_budget.pdf}
    \caption{the reaction budget}
  \end{subfigure}\hfill
  \begin{subfigure}[b]{0.44\linewidth}
    \centering\includegraphics[width=\linewidth]{gallery/baselines/posture_fallback_gap.pdf}
    \caption{the fallback assumption}
  \end{subfigure}
  \caption[The reaction budget and the fallback gap]{Left: the share of the
  \SI{150}{\milli\metre} safety gap a moving wearer closes before the robot
  knows, three limb-motion classes at the best (\SI{62}{\milli\second}) and
  worst (\SI{562}{\milli\second}) measured delay. Right: what the guard
  assumes when the wearer camera drops out (arms down, \SI{317}{\milli\metre})
  against a wearer with arms held out (\SI{93}{\milli\metre}).}
\end{figure}

\begin{figure}[H]
  \centering
  \includegraphics[width=\linewidth]{gallery/baselines/sim_real_park_error.pdf}
  \caption[The real arms park short by a fixed amount]{Thirty-six real moves
  of \SI{100}{\milli\metre}: every joint stops about \SI{0.29}{\degree} short
  (a), which costs the hand \SIrange{5}{8}{\milli\metre} by direction, the
  same at all three speeds (b).}
\end{figure}

\begin{figure}[H]
  \centering
  \includegraphics[width=\linewidth]{gallery/baselines/intent_estimator_sweep.pdf}
  \caption[Guessing the operator's intent from a typed-in map against a
  measured one]{Four hundred trials per point: from a typed-in object
  position the estimator's confident wrong guesses climb past the
  published \SI{16}{\percent} break-even at \SI{40}{\milli\metre} of error
  and reach \SI{35}{\percent} at one cube spacing; from the camera's
  measured position they stay at zero.}
\end{figure}

\begin{figure}[H]
  \centering
  \begin{subfigure}[b]{0.49\linewidth}
    \centering\includegraphics[width=\linewidth]{gallery/baselines/orientation_cost.pdf}
    \caption{what pinning the wrist costs}
  \end{subfigure}\hfill
  \begin{subfigure}[b]{0.49\linewidth}
    \centering\includegraphics[width=\linewidth]{gallery/baselines/what_binds_work_point.pdf}
    \caption{what stops the hand, by direction}
  \end{subfigure}
  \caption[The pinned wrist, and what stops the hand]{Left: mean reach under
  each wrist policy from home and from the work point. Right: from the work
  point, what stops the hand in each of twelve directions -- the wearer in
  eleven of twelve when the wrist is pinned.}
\end{figure}

\begin{figure}[H]
  \centering
  \begin{subfigure}[b]{0.49\linewidth}
    \centering\includegraphics[width=\linewidth]{gallery/baselines/wearer_posture.pdf}
    \caption{the wearer's posture}
  \end{subfigure}\hfill
  \begin{subfigure}[b]{0.49\linewidth}
    \centering\includegraphics[width=\linewidth]{gallery/baselines/reach_mount_sweep.pdf}
    \caption{moving the mounts}
  \end{subfigure}
  \caption[Wearer posture and mount position]{Left: how the wearer's arm
  posture changes what the arms can reach, with the \SI{150}{\milli\metre}
  gap drawn. Right: how moving the mounts changes what both arms can reach
  -- as built the two share three of 63 cells; the buildable best is 22.}
\end{figure}

\begin{figure}[H]
  \centering
  \begin{subfigure}[b]{0.49\linewidth}
    \centering\includegraphics[width=\linewidth]{gallery/baselines/pick_accuracy_belief_vs_miss.pdf}
    \caption{pad miss against camera error}
  \end{subfigure}\hfill
  \begin{subfigure}[b]{0.49\linewidth}
    \centering\includegraphics[width=\linewidth]{gallery/baselines/channel_health.pdf}
    \caption{the master's sensor channels}
  \end{subfigure}
  \caption[Pick accuracy, and channel health]{Left: whatever the camera
  believes, the arm delivers the fingers there -- the pad miss equals the
  belief error, so the residual is perception, not kinematics. Right: which
  of the fourteen master-arm sensors read cleanly, the verdict behind
  \cref{sec:channelhealth}.}
\end{figure}

\clearpage
\section{How the project was done}

\begin{figure}[H]
  \centering
  \includegraphics[width=\linewidth]{gallery/process/instrument_defects.pdf}
  \caption[The twenty-six instrument defects]{What went wrong in the
  twenty-six cases where a measurement was wrong and the system was not, and
  how each was caught: fifteen by an independently known quantity
  disagreeing with the reading, five by reading the code, six by looking at
  the raw image or data.}
\end{figure}

\begin{figure}[H]
  \centering
  \begin{subfigure}[b]{0.49\linewidth}
    \centering\includegraphics[width=\linewidth]{gallery/process/commits_and_tests.pdf}
    \caption{commits and tests}
  \end{subfigure}\hfill
  \begin{subfigure}[b]{0.49\linewidth}
    \centering\includegraphics[width=\linewidth]{gallery/process/recording_yield.pdf}
    \caption{recording yield}
  \end{subfigure}
  \caption[Commits, tests and recording yield]{Left: 449 commits over 27
  active days and a test suite that grew with the code to 140 files and
  1376 tests. Right: recording sessions started per day and how many held
  data -- unreliable until 2 September (0/6, 1/10, 4/17, 0/2, 0/9), then
  17/26, 3/3 and 14/15.}
\end{figure}
